% Options for packages loaded elsewhere
% Options for packages loaded elsewhere
\PassOptionsToPackage{unicode}{hyperref}
\PassOptionsToPackage{hyphens}{url}
\PassOptionsToPackage{dvipsnames,svgnames,x11names}{xcolor}
%
\documentclass[
  letterpaper,
  DIV=11,
  numbers=noendperiod]{scrartcl}
\usepackage{xcolor}
\usepackage{amsmath,amssymb}
\setcounter{secnumdepth}{-\maxdimen} % remove section numbering
\usepackage{iftex}
\ifPDFTeX
  \usepackage[T1]{fontenc}
  \usepackage[utf8]{inputenc}
  \usepackage{textcomp} % provide euro and other symbols
\else % if luatex or xetex
  \usepackage{unicode-math} % this also loads fontspec
  \defaultfontfeatures{Scale=MatchLowercase}
  \defaultfontfeatures[\rmfamily]{Ligatures=TeX,Scale=1}
\fi
\usepackage{lmodern}
\ifPDFTeX\else
  % xetex/luatex font selection
\fi
% Use upquote if available, for straight quotes in verbatim environments
\IfFileExists{upquote.sty}{\usepackage{upquote}}{}
\IfFileExists{microtype.sty}{% use microtype if available
  \usepackage[]{microtype}
  \UseMicrotypeSet[protrusion]{basicmath} % disable protrusion for tt fonts
}{}
\makeatletter
\@ifundefined{KOMAClassName}{% if non-KOMA class
  \IfFileExists{parskip.sty}{%
    \usepackage{parskip}
  }{% else
    \setlength{\parindent}{0pt}
    \setlength{\parskip}{6pt plus 2pt minus 1pt}}
}{% if KOMA class
  \KOMAoptions{parskip=half}}
\makeatother
% Make \paragraph and \subparagraph free-standing
\makeatletter
\ifx\paragraph\undefined\else
  \let\oldparagraph\paragraph
  \renewcommand{\paragraph}{
    \@ifstar
      \xxxParagraphStar
      \xxxParagraphNoStar
  }
  \newcommand{\xxxParagraphStar}[1]{\oldparagraph*{#1}\mbox{}}
  \newcommand{\xxxParagraphNoStar}[1]{\oldparagraph{#1}\mbox{}}
\fi
\ifx\subparagraph\undefined\else
  \let\oldsubparagraph\subparagraph
  \renewcommand{\subparagraph}{
    \@ifstar
      \xxxSubParagraphStar
      \xxxSubParagraphNoStar
  }
  \newcommand{\xxxSubParagraphStar}[1]{\oldsubparagraph*{#1}\mbox{}}
  \newcommand{\xxxSubParagraphNoStar}[1]{\oldsubparagraph{#1}\mbox{}}
\fi
\makeatother

\usepackage{color}
\usepackage{fancyvrb}
\newcommand{\VerbBar}{|}
\newcommand{\VERB}{\Verb[commandchars=\\\{\}]}
\DefineVerbatimEnvironment{Highlighting}{Verbatim}{commandchars=\\\{\}}
% Add ',fontsize=\small' for more characters per line
\usepackage{framed}
\definecolor{shadecolor}{RGB}{241,243,245}
\newenvironment{Shaded}{\begin{snugshade}}{\end{snugshade}}
\newcommand{\AlertTok}[1]{\textcolor[rgb]{0.68,0.00,0.00}{#1}}
\newcommand{\AnnotationTok}[1]{\textcolor[rgb]{0.37,0.37,0.37}{#1}}
\newcommand{\AttributeTok}[1]{\textcolor[rgb]{0.40,0.45,0.13}{#1}}
\newcommand{\BaseNTok}[1]{\textcolor[rgb]{0.68,0.00,0.00}{#1}}
\newcommand{\BuiltInTok}[1]{\textcolor[rgb]{0.00,0.23,0.31}{#1}}
\newcommand{\CharTok}[1]{\textcolor[rgb]{0.13,0.47,0.30}{#1}}
\newcommand{\CommentTok}[1]{\textcolor[rgb]{0.37,0.37,0.37}{#1}}
\newcommand{\CommentVarTok}[1]{\textcolor[rgb]{0.37,0.37,0.37}{\textit{#1}}}
\newcommand{\ConstantTok}[1]{\textcolor[rgb]{0.56,0.35,0.01}{#1}}
\newcommand{\ControlFlowTok}[1]{\textcolor[rgb]{0.00,0.23,0.31}{\textbf{#1}}}
\newcommand{\DataTypeTok}[1]{\textcolor[rgb]{0.68,0.00,0.00}{#1}}
\newcommand{\DecValTok}[1]{\textcolor[rgb]{0.68,0.00,0.00}{#1}}
\newcommand{\DocumentationTok}[1]{\textcolor[rgb]{0.37,0.37,0.37}{\textit{#1}}}
\newcommand{\ErrorTok}[1]{\textcolor[rgb]{0.68,0.00,0.00}{#1}}
\newcommand{\ExtensionTok}[1]{\textcolor[rgb]{0.00,0.23,0.31}{#1}}
\newcommand{\FloatTok}[1]{\textcolor[rgb]{0.68,0.00,0.00}{#1}}
\newcommand{\FunctionTok}[1]{\textcolor[rgb]{0.28,0.35,0.67}{#1}}
\newcommand{\ImportTok}[1]{\textcolor[rgb]{0.00,0.46,0.62}{#1}}
\newcommand{\InformationTok}[1]{\textcolor[rgb]{0.37,0.37,0.37}{#1}}
\newcommand{\KeywordTok}[1]{\textcolor[rgb]{0.00,0.23,0.31}{\textbf{#1}}}
\newcommand{\NormalTok}[1]{\textcolor[rgb]{0.00,0.23,0.31}{#1}}
\newcommand{\OperatorTok}[1]{\textcolor[rgb]{0.37,0.37,0.37}{#1}}
\newcommand{\OtherTok}[1]{\textcolor[rgb]{0.00,0.23,0.31}{#1}}
\newcommand{\PreprocessorTok}[1]{\textcolor[rgb]{0.68,0.00,0.00}{#1}}
\newcommand{\RegionMarkerTok}[1]{\textcolor[rgb]{0.00,0.23,0.31}{#1}}
\newcommand{\SpecialCharTok}[1]{\textcolor[rgb]{0.37,0.37,0.37}{#1}}
\newcommand{\SpecialStringTok}[1]{\textcolor[rgb]{0.13,0.47,0.30}{#1}}
\newcommand{\StringTok}[1]{\textcolor[rgb]{0.13,0.47,0.30}{#1}}
\newcommand{\VariableTok}[1]{\textcolor[rgb]{0.07,0.07,0.07}{#1}}
\newcommand{\VerbatimStringTok}[1]{\textcolor[rgb]{0.13,0.47,0.30}{#1}}
\newcommand{\WarningTok}[1]{\textcolor[rgb]{0.37,0.37,0.37}{\textit{#1}}}

\usepackage{longtable,booktabs,array}
\newcounter{none} % for unnumbered tables
\usepackage{calc} % for calculating minipage widths
% Correct order of tables after \paragraph or \subparagraph
\usepackage{etoolbox}
\makeatletter
\patchcmd\longtable{\par}{\if@noskipsec\mbox{}\fi\par}{}{}
\makeatother
% Allow footnotes in longtable head/foot
\IfFileExists{footnotehyper.sty}{\usepackage{footnotehyper}}{\usepackage{footnote}}
\makesavenoteenv{longtable}
\usepackage{graphicx}
\makeatletter
\newsavebox\pandoc@box
\newcommand*\pandocbounded[1]{% scales image to fit in text height/width
  \sbox\pandoc@box{#1}%
  \Gscale@div\@tempa{\textheight}{\dimexpr\ht\pandoc@box+\dp\pandoc@box\relax}%
  \Gscale@div\@tempb{\linewidth}{\wd\pandoc@box}%
  \ifdim\@tempb\p@<\@tempa\p@\let\@tempa\@tempb\fi% select the smaller of both
  \ifdim\@tempa\p@<\p@\scalebox{\@tempa}{\usebox\pandoc@box}%
  \else\usebox{\pandoc@box}%
  \fi%
}
% Set default figure placement to htbp
\def\fps@figure{htbp}
\makeatother





\setlength{\emergencystretch}{3em} % prevent overfull lines

\providecommand{\tightlist}{%
  \setlength{\itemsep}{0pt}\setlength{\parskip}{0pt}}



 


\KOMAoption{captions}{tableheading}
\makeatletter
\@ifpackageloaded{caption}{}{\usepackage{caption}}
\AtBeginDocument{%
\ifdefined\contentsname
  \renewcommand*\contentsname{Table of contents}
\else
  \newcommand\contentsname{Table of contents}
\fi
\ifdefined\listfigurename
  \renewcommand*\listfigurename{List of Figures}
\else
  \newcommand\listfigurename{List of Figures}
\fi
\ifdefined\listtablename
  \renewcommand*\listtablename{List of Tables}
\else
  \newcommand\listtablename{List of Tables}
\fi
\ifdefined\figurename
  \renewcommand*\figurename{Figure}
\else
  \newcommand\figurename{Figure}
\fi
\ifdefined\tablename
  \renewcommand*\tablename{Table}
\else
  \newcommand\tablename{Table}
\fi
}
\@ifpackageloaded{float}{}{\usepackage{float}}
\floatstyle{ruled}
\@ifundefined{c@chapter}{\newfloat{codelisting}{h}{lop}}{\newfloat{codelisting}{h}{lop}[chapter]}
\floatname{codelisting}{Listing}
\newcommand*\listoflistings{\listof{codelisting}{List of Listings}}
\makeatother
\makeatletter
\makeatother
\makeatletter
\@ifpackageloaded{caption}{}{\usepackage{caption}}
\@ifpackageloaded{subcaption}{}{\usepackage{subcaption}}
\makeatother
\usepackage{bookmark}
\IfFileExists{xurl.sty}{\usepackage{xurl}}{} % add URL line breaks if available
\urlstyle{same}
\hypersetup{
  pdftitle={Rating factors shouldn't depend on who's in the book},
  pdfauthor={Ariel Balter},
  colorlinks=true,
  linkcolor={blue},
  filecolor={Maroon},
  citecolor={Blue},
  urlcolor={Blue},
  pdfcreator={LaTeX via pandoc}}


\title{Rating factors shouldn't depend on who's in the book}
\usepackage{etoolbox}
\makeatletter
\providecommand{\subtitle}[1]{% add subtitle to \maketitle
  \apptocmd{\@title}{\par {\large #1 \par}}{}{}
}
\makeatother
\subtitle{Relativities, coefficients, and the interaction decomposition}
\author{Ariel Balter}
\date{October 23, 2026}
\begin{document}
\maketitle


\subsection{The problem}\label{problem}

\textbf{``Relativity'' names at least four different numbers.}

{\def\LTcaptype{none} % do not increment counter
\begin{longtable}[]{@{}
  >{\raggedright\arraybackslash}p{(\linewidth - 2\tabcolsep) * \real{0.5000}}
  >{\raggedright\arraybackslash}p{(\linewidth - 2\tabcolsep) * \real{0.5000}}@{}}
\toprule\noalign{}
\begin{minipage}[b]{\linewidth}\raggedright
what people call ``the relativity''
\end{minipage} & \begin{minipage}[b]{\linewidth}\raggedright
how it's computed
\end{minipage} \\
\midrule\noalign{}
\endhead
\bottomrule\noalign{}
\endlastfoot
one-way & average loss cost at a level \(\div\) at the base level \\
GLM factor & \(e^{\beta}\) relative to a base level \\
partial dependence & hold the variable fixed, average the model over the
book \\
model one-way & average model prediction by level \\
\end{longtable}
}

\textbf{Only one is a property of the model alone.} The others are
averages over the book --- they move when the book moves, with the model
unchanged.

Why it matters here:

\begin{itemize}
\tightlist
\item
  Our GAM has no table to read off. Something has to be extracted, and
  the default extraction (partial dependence) is a book average.
\item
  AvT: a mix shift reads as mispricing.
\item
  Factors aren't comparable across time, segments or filings unless the
  reference book travels with them.
\end{itemize}

Open with the sentence on the title slide --- that's the whole talk. A
rating factor is supposed to say what the model thinks about a breed, an
age, a ZIP. If the number changes because the mix of other pets in the
book changed, it isn't telling you about the breed; it's telling you
about the book.

Walk the four rows. The point of the table is not that these are wrong
--- each is a legitimate quantity --- but that we use one word for all
of them and only one of them (the GLM factor) is intrinsic to the model.

Why now: we're pricing off a GAM. There is no coefficient table in the
sense a GLM has one; mgcv's coefficients are centered on the training
data by identifiability constraints, so even ``read off the
coefficients'' is a book-dependent gauge until you re-gauge it. The
natural thing to reach for is a partial-dependence curve, and that is
what the ML-in-insurance literature does when it turns a GBM into a
tariff.

If asked ``we already know one-ways are distorted by correlation'': yes
--- that's rung one on the next slide, it's the textbook argument for
GLMs. The partial-dependence curve was supposed to fix that, and it
fixes half of it.

If asked ``we use GLM coefficients, not PDPs'': yes, and that's the
right answer. The paper proves why, and gives the recipe for getting the
same thing out of a GAM or any other model --- which is what we need.

\subsection{Three ways to extract a factor}\label{ladder}

Each rung fixes the failure of the one below it.

{\def\LTcaptype{none} % do not increment counter
\begin{longtable}[]{@{}
  >{\raggedright\arraybackslash}p{(\linewidth - 6\tabcolsep) * \real{0.1667}}
  >{\centering\arraybackslash}p{(\linewidth - 6\tabcolsep) * \real{0.2778}}
  >{\centering\arraybackslash}p{(\linewidth - 6\tabcolsep) * \real{0.2778}}
  >{\centering\arraybackslash}p{(\linewidth - 6\tabcolsep) * \real{0.2778}}@{}}
\toprule\noalign{}
\begin{minipage}[b]{\linewidth}\raggedright
\end{minipage} & \begin{minipage}[b]{\linewidth}\centering
correlation with other rating variables
\end{minipage} & \begin{minipage}[b]{\linewidth}\centering
mix of interacting partners
\end{minipage} & \begin{minipage}[b]{\linewidth}\centering
product of tables reproduces the model
\end{minipage} \\
\midrule\noalign{}
\endhead
\bottomrule\noalign{}
\endlastfoot
\textbf{one-way} & leaks in & leaks in & no \\
\textbf{partial dependence} (PDP) & removed & leaks in whenever an
interaction term straddles the variable & no --- the pair table
double-counts the main effects \\
\textbf{decomposition anchored at the base cell} & removed & removed &
yes, exactly \\
\end{longtable}
}

The third row is what GLM treatment coding does automatically. The work
is to do it \emph{on purpose}, for any model, and to know \emph{when}
the second row is good enough.

Rung one: one-way. Averages over everything, so if older pets are
disproportionately one breed, the breed factor absorbs age. Everyone in
the room knows this; it's why we fit GLMs.

Rung two: partial dependence. Set the breed for every row in the book,
average the predictions, compare levels. Correlation is gone because
every row gets every breed. But if the model has a breed-by-age
interaction, the breed curve is averaged over the book's age mix, so the
breed factor still depends on how many old pets we have. The precise
statement: the partial relativity for a variable set S is intrinsic to
the model if and only if no model term partially overlaps S. A term that
straddles S leaves a book-dependent piece behind.

Rung three: the decomposition. Build a table for every subset of
variables that a model term touches, and make each table hold only the
part of the effect that belongs to that subset --- subtract what belongs
to its sub-tables (inclusion--exclusion). If ``subtract'' means
``compare to the base cell,'' the tables are exactly treatment-coded
coefficients: mix-free, 1.00 at base, and their product reproduces the
model to the penny. If ``subtract'' means ``average over the book,'' the
product still reproduces the model but the split between main effect and
interaction depends on the mix.

The ``double-counts'' cell: if you build PDP tables for breed, age and
breed-by-age and multiply, the pair table already contains both main
effects. GLM people never see this because treatment coding has already
done the subtraction.

If asked ``isn't this just functional ANOVA?'': yes, and the paper says
so and cites it. The contribution is the term-level condition for when a
table is intrinsic, the base-cell-versus-book-average choice framed as a
rating-table decision, and one vocabulary that covers all four rows of
the previous slide.

One caveat worth saying out loud: ``mix-free'' means mix-free given the
model. If the model is missing an interaction, the fitted main effect
absorbed it in a mix-dependent way, and no extraction fixes that.
Modeling and factor reporting have to be done together.

\subsection{Worked example: nothing about the model
changes}\label{example}

Two breeds \(\times\) two ages, log link, base cell = Breed A / young.
Breed B \(\times 1.5\), old \(\times 1.2\), B-and-old interaction
\(\times 1.5\). So the four cells are A/young 1.00, A/old 1.20, B/young
1.50, B/old 2.70.

{\def\LTcaptype{none} % do not increment counter
\begin{longtable}[]{@{}
  >{\raggedright\arraybackslash}p{(\linewidth - 6\tabcolsep) * \real{0.1667}}
  >{\centering\arraybackslash}p{(\linewidth - 6\tabcolsep) * \real{0.2778}}
  >{\centering\arraybackslash}p{(\linewidth - 6\tabcolsep) * \real{0.2778}}
  >{\centering\arraybackslash}p{(\linewidth - 6\tabcolsep) * \real{0.2778}}@{}}
\toprule\noalign{}
\begin{minipage}[b]{\linewidth}\raggedright
book (share old)
\end{minipage} & \begin{minipage}[b]{\linewidth}\centering
one-way breed factor
\end{minipage} & \begin{minipage}[b]{\linewidth}\centering
PDP breed factor
\end{minipage} & \begin{minipage}[b]{\linewidth}\centering
coefficient
\end{minipage} \\
\midrule\noalign{}
\endhead
\bottomrule\noalign{}
\endlastfoot
20 \% & 1.67 & 1.63 & \textbf{1.50} \\
50 \% & 1.91 & 1.84 & \textbf{1.50} \\
80 \% & 2.12 & 2.07 & \textbf{1.50} \\
50 \% overall, but B skews old (A 30 \%, B 70 \%) & 2.21 & 1.84 &
\textbf{1.50} \\
\end{longtable}
}

The PDP column has a closed form --- and it \emph{is} the theorem for
two variables:

\[
R^{\text{PDP}}_{B} \;=\; \exp\!\big(\beta_B + p_{\text{old}}\,\beta_{B\times\text{old}}\big) \;=\; 1.5 \times 1.5^{\,p_{\text{old}}}
\]

\textbf{The partial relativity is the coefficient, plus the interaction
coefficient times the share of old pets in the book.}

Everything on this slide uses one fixed model. Only the book changes
between rows.

Read down the one-way column: it moves with the age mix (rows 1--3) and
it moves with breed--age correlation even at a fixed overall mix (row
4). Read down the PDP column: row 4 equals row 2 --- correlation is gone
--- but rows 1--3 still move, because the breed-by-age term straddles
``breed.'' The coefficient column doesn't move.

The formula is the slide's payoff. Beta-B is the breed coefficient,
beta-B-times-old is the interaction, p-old is the share of old pets in
whatever book you averaged over. In this toy, 1.5 to the power one plus
p-old: from 1.50 if the book is all young to 2.25 if it's all old, for
the same model. Say the bold sentence slowly.

Connect to AvT: if the age mix of a segment drifts, its ``breed
relativity'' drifts, and actual-versus-technical reads it as mispricing.
Nothing was mispriced.

One-way here is the ratio of mean loss cost by breed; PDP is on the link
scale, exponentiated. If someone asks about the scale: that's the
convention for a multiplicative plan, and the paper is explicit about
it.

The real breed-by-age surface from our GAM is the version that goes in
the guide. This toy is just small enough to see.

\subsection{Two deliverables}\label{deliverables}

\subsubsection{Practical guide
(internal)}\label{practical-guide-internal}

\begin{itemize}
\tightlist
\item
  One vocabulary: which of the four numbers you mean, and when each is
  the right one
\item
  The algebra as three operators (\texttt{filter}, \texttt{mutate},
  weighted mean) --- every relativity is a two-line \texttt{dplyr} pipe
\item
  The recommendation, with the proof: rating factors = base-cell term
  contributions, for GLM or GAM
\item
  The breed \(\times\) age demonstration from our model
\end{itemize}

\subsubsection{Journal article
(external)}\label{journal-article-external}

\begin{itemize}
\tightlist
\item
  Same framework, with the theorem and proofs
\item
  The term-level condition for when a relativity is intrinsic
  vs.~portfolio-defined
\item
  Demonstrations on public insurance datasets, comparable to the
  published GBM-to-tariff work
\item
  Peer review as quality control on a method we're about to rate on
\end{itemize}

From the paper's one-page summary:

\begin{itemize}
\tightlist
\item
  A one-way partial dependence plot is generally \textbf{not} ``the
  relativity'' for a variable.
\item
  Portfolio mix can change reported marginal relativities while the
  fitted coefficients are unchanged.
\item
  Rating manuals should be built from coefficient and interaction terms,
  not portfolio-averaged effects.
\item
  Reference portfolios should be disclosed whenever marginal
  relativities are reported.
\end{itemize}

The guide is the thing colleagues use next quarter. It answers three
questions: what do you mean by relativity, how do you compute the one
you mean, and which one goes in the rating table. The operator-to-dplyr
mapping is the part people like most --- the math is three operations
and each one is a verb they already use.

The article is the same material with the proofs and public-data
demonstrations. Framing for the room: we're about to rate off a method;
having it refereed is cheap insurance, and it's the kind of thing that's
good for the company's name in the actuarial literature.

The four bullets are lifted from the one-page summary, which is in the
pre-read.

\subsection{The ask}\label{ask}

\textbf{Three weeks, about 15 hours a week --- 45 hours.}

\subsubsection{Done}\label{done}

\begin{itemize}
\tightlist
\item
  Unified vocabulary and algebra
\item
  Operator notation → the two-line code it produces
\item
  The decomposition proof
\item
  Paper sections drafted
\item
  One-page summary (pre-read)
\end{itemize}

\subsubsection{Remaining}\label{remaining}

\begin{itemize}
\tightlist
\item
  More worked examples, including breed \(\times\) age from our model
\item
  Demonstrations on open-source insurance data (French MTPL, Belgian
  MTPL, Australian auto)
\item
  Guide: assemble, review with the team
\item
  Article: polish, format, submit
\end{itemize}

\textbf{Sequence:} guide first (weeks 1--2, immediately usable, review
checkpoint) --- article last (week 3).

Be concrete about the hours and about what exists already --- most of
the thinking is done; what's left is examples, assembly and polish.

The public datasets are the ones the published GBM-to-tariff papers use
(freMTPL2 via CASdatasets, the Belgian MTPL data, ausprivauto0405), so
our results are directly comparable. There is no public pet-insurance
claims dataset I know of; the pet demonstration stays internal, in the
guide.

The guide-first sequencing is the risk control: by the end of week two
there is a usable internal document regardless of what happens with the
article.

Close by restating the title: the value is that our rating factors stop
depending on who happens to be in the book.

\section*{Backup}\label{backup}
\addcontentsline{toc}{section}{Backup}

\subsection{The result in one line}\label{theorem}

Link-scale predictor as a sum of terms, each on a subset \(t\) of the
variables:

\[
F(x) \;=\; \sum_{t \in T} f_t(x_t)
\]

Partial relativity for a variable set \(S\) --- set \(x_S\) to \(z_S\)
on every row, take the weighted mean over the book \(w\):

\[
M_S(z_S) \;=\; \mathbb{E}_w\big[\,F(z_S, X_{-S})\,\big]
\]

\textbf{\(M_S\) is intrinsic to the model iff every term satisfies
\(t \subseteq S\) or \(t \cap S = \varnothing\).} A straddling term
contributes \(\mathbb{E}_w\, f_t(z_{t \cap S},\, X_{t \setminus S})\)
--- a function of the book.

Inclusion--exclusion gives the \(S\)-specific part --- the same formula
for any ``remove'' rule:

\[
g_S(z_S) \;=\; \sum_{T \subseteq S} (-1)^{|S \setminus T|}\, M_T(z_T)
\]

\begin{itemize}
\tightlist
\item
  Remove = \emph{average over the book}: product reproduces the model;
  the main-effect / interaction split depends on the mix.
\item
  Remove = \emph{set to the base cell}, \(M_T(z_T) = F(z_T, x^0_{-T})\):
  product reproduces the model, mix-free, and the \(g_S\) are exactly
  the treatment-coded coefficients.
\item
  \(g_S = 0\) unless \(S\) lies inside some model term --- only subsets
  of actual terms need tables.
\end{itemize}

Only if someone wants the mathematics. The two-variable case of the
inclusion--exclusion formula is the familiar one: g for the pair equals
the pair table minus the two single tables plus the constant.

The last bullet is the practical safeguard: the number of tables does
not explode; you need tables for the subsets of terms that are actually
in the model.

\subsection{The code behind the example}\label{code}

\begin{Shaded}
\begin{Highlighting}[]
\FunctionTok{library}\NormalTok{(dplyr)}

\NormalTok{cells }\OtherTok{\textless{}{-}} \FunctionTok{tibble}\NormalTok{(}\AttributeTok{breed =} \FunctionTok{c}\NormalTok{(}\StringTok{"A"}\NormalTok{, }\StringTok{"A"}\NormalTok{, }\StringTok{"B"}\NormalTok{, }\StringTok{"B"}\NormalTok{),}
                \AttributeTok{age   =} \FunctionTok{c}\NormalTok{(}\StringTok{"young"}\NormalTok{, }\StringTok{"old"}\NormalTok{, }\StringTok{"young"}\NormalTok{, }\StringTok{"old"}\NormalTok{),}
                \AttributeTok{mu    =} \FunctionTok{c}\NormalTok{(}\FloatTok{1.00}\NormalTok{, }\FloatTok{1.20}\NormalTok{, }\FloatTok{1.50}\NormalTok{, }\FloatTok{2.70}\NormalTok{))}

\NormalTok{book }\OtherTok{\textless{}{-}} \ControlFlowTok{function}\NormalTok{(p\_old\_A, p\_old\_B)}
  \FunctionTok{tibble}\NormalTok{(}\AttributeTok{breed =} \FunctionTok{c}\NormalTok{(}\StringTok{"A"}\NormalTok{, }\StringTok{"A"}\NormalTok{, }\StringTok{"B"}\NormalTok{, }\StringTok{"B"}\NormalTok{),}
         \AttributeTok{age   =} \FunctionTok{c}\NormalTok{(}\StringTok{"young"}\NormalTok{, }\StringTok{"old"}\NormalTok{, }\StringTok{"young"}\NormalTok{, }\StringTok{"old"}\NormalTok{),}
         \AttributeTok{w     =} \FunctionTok{c}\NormalTok{(}\DecValTok{1} \SpecialCharTok{{-}}\NormalTok{ p\_old\_A, p\_old\_A, }\DecValTok{1} \SpecialCharTok{{-}}\NormalTok{ p\_old\_B, p\_old\_B) }\SpecialCharTok{/} \DecValTok{2}\NormalTok{)}

\CommentTok{\# one{-}way: weighted mean loss cost by breed, B over A}
\NormalTok{one\_way }\OtherTok{\textless{}{-}} \ControlFlowTok{function}\NormalTok{(bk) bk }\SpecialCharTok{|\textgreater{}}
  \FunctionTok{inner\_join}\NormalTok{(cells, }\AttributeTok{by =} \FunctionTok{c}\NormalTok{(}\StringTok{"breed"}\NormalTok{, }\StringTok{"age"}\NormalTok{)) }\SpecialCharTok{|\textgreater{}}
  \FunctionTok{group\_by}\NormalTok{(breed) }\SpecialCharTok{|\textgreater{}} \FunctionTok{summarise}\NormalTok{(}\AttributeTok{m =} \FunctionTok{sum}\NormalTok{(w }\SpecialCharTok{*}\NormalTok{ mu) }\SpecialCharTok{/} \FunctionTok{sum}\NormalTok{(w)) }\SpecialCharTok{|\textgreater{}}
  \FunctionTok{with}\NormalTok{(m[breed }\SpecialCharTok{==} \StringTok{"B"}\NormalTok{] }\SpecialCharTok{/}\NormalTok{ m[breed }\SpecialCharTok{==} \StringTok{"A"}\NormalTok{])}

\CommentTok{\# partial: set breed on every row (mutate), weighted mean of log mu}
\NormalTok{partial }\OtherTok{\textless{}{-}} \ControlFlowTok{function}\NormalTok{(bk, z) bk }\SpecialCharTok{|\textgreater{}}
  \FunctionTok{mutate}\NormalTok{(}\AttributeTok{breed =}\NormalTok{ z) }\SpecialCharTok{|\textgreater{}}
  \FunctionTok{inner\_join}\NormalTok{(cells, }\AttributeTok{by =} \FunctionTok{c}\NormalTok{(}\StringTok{"breed"}\NormalTok{, }\StringTok{"age"}\NormalTok{)) }\SpecialCharTok{|\textgreater{}}
  \FunctionTok{summarise}\NormalTok{(}\AttributeTok{m =} \FunctionTok{sum}\NormalTok{(w }\SpecialCharTok{*} \FunctionTok{log}\NormalTok{(mu)) }\SpecialCharTok{/} \FunctionTok{sum}\NormalTok{(w)) }\SpecialCharTok{|\textgreater{}} \FunctionTok{pull}\NormalTok{(m)}

\NormalTok{pdp\_rel }\OtherTok{\textless{}{-}} \ControlFlowTok{function}\NormalTok{(bk) }\FunctionTok{exp}\NormalTok{(}\FunctionTok{partial}\NormalTok{(bk, }\StringTok{"B"}\NormalTok{) }\SpecialCharTok{{-}} \FunctionTok{partial}\NormalTok{(bk, }\StringTok{"A"}\NormalTok{))}

\CommentTok{\# coefficient: everything at the base cell, then breed to B}
\NormalTok{coef\_rel }\OtherTok{\textless{}{-}} \FunctionTok{with}\NormalTok{(cells, mu[breed }\SpecialCharTok{==} \StringTok{"B"} \SpecialCharTok{\&}\NormalTok{ age }\SpecialCharTok{==} \StringTok{"young"}\NormalTok{] }\SpecialCharTok{/}
\NormalTok{                        mu[breed }\SpecialCharTok{==} \StringTok{"A"} \SpecialCharTok{\&}\NormalTok{ age }\SpecialCharTok{==} \StringTok{"young"}\NormalTok{])}

\NormalTok{books }\OtherTok{\textless{}{-}} \FunctionTok{list}\NormalTok{(}\FunctionTok{book}\NormalTok{(.}\DecValTok{2}\NormalTok{, .}\DecValTok{2}\NormalTok{), }\FunctionTok{book}\NormalTok{(.}\DecValTok{5}\NormalTok{, .}\DecValTok{5}\NormalTok{), }\FunctionTok{book}\NormalTok{(.}\DecValTok{8}\NormalTok{, .}\DecValTok{8}\NormalTok{), }\FunctionTok{book}\NormalTok{(.}\DecValTok{3}\NormalTok{, .}\DecValTok{7}\NormalTok{))}
\FunctionTok{sapply}\NormalTok{(books, one\_way)   }\CommentTok{\# 1.673 1.909 2.121 2.208}
\FunctionTok{sapply}\NormalTok{(books, pdp\_rel)   }\CommentTok{\# 1.627 1.837 2.075 1.837}
\NormalTok{coef\_rel                 }\CommentTok{\# 1.5}
\end{Highlighting}
\end{Shaded}

This is the ``every relativity is a two-line pipe'' claim made good:
one-way is a group-by, partial is a mutate followed by a weighted mean,
the coefficient is a mutate of everything to the base row.

\subsection{Sources and prior art}\label{sources}

\textbf{Current practice}

\begin{itemize}
\tightlist
\item
  Werner \& Modlin, \emph{Basic Ratemaking} (CAS) --- one-way
  vs.~multivariate classification
\item
  Anderson et al.~2007, \emph{A Practitioner's Guide to GLMs} (CAS);
  Goldburd, Khare, Tevet \& Guller, \emph{GLMs for Insurance Rating}
  (CAS Monograph 5) --- coefficients to rating factors, interactions,
  offsets
\item
  Henckaerts, Côté, Antonio \& Verbelen 2021 (NAAJ); Henckaerts, Antonio
  \& Côté 2022 (maidrr) --- tariffs from GBMs via partial dependence
\item
  Lorentzen \& Mayer 2020; Wüthrich \& Merz 2023 --- PDP/ALE as the
  interpretation tools in actuarial ML
\end{itemize}

\textbf{The mathematics}

\begin{itemize}
\tightlist
\item
  Friedman 2001 --- partial dependence defined as an average over the
  book
\item
  Hooker 2007; Apley \& Zhu 2020 --- PDP equals the functional-ANOVA
  main effect only under independence
\item
  Kuo, Sloan, Wasilkowski \& Woźniakowski 2010; Rabitz \& Alış 1999 ---
  anchored vs.~ANOVA decompositions (base cell vs.~book average)
\item
  Lengerich et al.~2020; Herren \& Hahn 2022; Hiabu, Meyer \& Wright
  2023; Bordt \& von Luxburg 2023 --- purified interactions, Möbius
  inversion, Shapley connections
\end{itemize}

Verify years and venues against your copies before this goes anywhere.
The ``current practice'' block is what backs slide 1; the
``mathematics'' block is what a referee will expect the article to
position against.




\end{document}
